\documentclass[book.tex]{subfiles}

\begin{document}
\chapter{Wick's Theorem and Feynman Diagrams}

\label{chap:feynman_diagrams_more}
\noindent\rule{11cm}{0.4pt} \\
\textbf{Prerequisites:} \autoref{chap:qft_basics}  \\
\textbf{Difficulty Level:} ** \\
\noindent\rule{11cm}{0.4pt}
\vspace{5mm}

%\textcolor{red}{TODO: Cite Peskin and Schroeder}

Feynman diagrams arise as a tool to organize complex perturbative calculations arising in quantum mechanics and quantum field theory. In particular, we can view Feynman diagrams as indexing terms in the Dyson series.

Wick's theorem provides the mathematical basis for the perturbative expansion that underpins Feynman diagrams.

\section{Normal Ordering}

A product of quantum fields is said to be normal ordered if all creation operators are to the left of all annihilation operators. If a product of operators is denoted by $O$, the normal-ordering of this product is denoted $:O:$. To give a few examples,
\begin{align}
    :\hat{a}^\dagger \hat{a}: &= \hat{a}^\dagger \hat{a}\\
    :\hat{a} \hat{a}^\dagger: &= \hat{a}^\dagger \hat{a}\\    
\end{align}
where $\hat{a}^\dagger$ is a creation operator and $\hat{a}$ is an annihilation operator.
\section{Wick's Theorem}

Wick's theorem is a tool for simplifying the computation of a higher order derivative into a combinatorial sum of lower order terms.

%\begin{align}
%    \contraction{}{A}{B}{C} \contraction[2ex]{A}{B}{C}{D} ABCD
%\end{align}

We define the contraction of an operator to be the difference between the product and the normal form of the product %(\textcolor{red}{TODO: Define the normal order})
\begin{align}
\contraction{}{A}{}{B} AB &= AB - :AB:
\end{align}

The general statement is given by
\begin{align}
    ABCDEF\dotsc &= :ABCDEF \dotsc : + \sum_{\textrm{single contractions}} :\contraction{}{A}{}{B} ABCDEF\dotsc : \
    &+ \sum_{\textrm{double contractions}} :\contraction{}{A}{}{B}
    \contraction{A}{B}{}{C}
    ABCDEF\dotsc : + \dotsc 
\end{align}
\subsection{Proof Sketch}
Let us consider the $N=2$ case of Wick's theorem. In this case, the formula simplifies to
\begin{align}
    AB &= :AB: + :\contraction{}{A}{}{B} AB:
\end{align}
This statement is true by definition. To prove the more general form, we proceed by induction. The new operator we add for the $N$ operator case must either be a creation or annihiliation operator. We can simplify on these two cases using the inductive assumption.

\section{An Alternative Formulation}

Suppose that we have $N$ real random variables that form a vector 

\begin{align}
    |X\rangle &= [x_1,\dotsc,x_n],\quad x_i \in \mathbb{R}
\end{align}
We define a partition function over these variables by
\begin{align}
    Z(J) &= \prod_{i=1}^N \int_{-\infty}^{\infty} dx_i e^{-S_0(X) + \langle J | X \rangle}
\end{align}
Here $J$ is a source field vector
\begin{align}
    |J\rangle &= [J_1, \dotsc, J_n] \\
    \langle J | X \rangle &= \sum_{i=1}^N J_i x_i
\end{align}
Recall that for a $1$-dimensional example, we have the Gaussian integral
\begin{align}
    \int_{-\infty}^{\infty} e^{-\frac{1}{2}mx^2 + Jx} &= \sqrt{\frac{2\pi}{m}} \exp \Bigg [ \frac{J^2}{2m} \Bigg ]
\end{align}
Let's now consider the matrix action
\begin{align}
    S(J) &= -\frac{1}{2}\langle X | M | X \rangle + \langle J | X \rangle
\end{align}
Here let $M$ be a real symmetric matrix that has orthonormal eigen vectors
\begin{align}
    M|\alpha \rangle &= m_{\alpha}|\alpha \rangle \\
    \langle \alpha | \beta \rangle &= \delta_{ab}
\end{align}
We can expand vectors $|X\rangle$ and $|J\rangle$ in terms of this orthonormal basis
\begin{align}
    |X\rangle &= \sum_{\alpha}x_\alpha |\alpha \rangle \\
    |J\rangle &= \sum_{\alpha}J_\alpha |\alpha \rangle \\
\end{align}
These two equations mean that
\begin{align}
    \langle J | X \rangle &= \sum_{\alpha}J_\alpha x_\alpha
\end{align}

Since the mapping $x_i \mapsto x_\alpha$ is orthogonal, the jacobian $\mathcal{J}$ is identity. We can rewrite the action $S$ now as
\begin{align}
    S(J) &= -\frac{1}{2}\langle X | M | X \rangle + \langle J | X \rangle \\
    &= \sum_{\alpha} \left ( -\frac{1}{2}x_\alpha^2m_\alpha + J_\alpha x_{\alpha} \right )
\end{align}
We can now factor the partition function into a producct of Gaussian integrals
\begin{align}
    Z(J) &= \prod_{i=1}^N \int_{-\infty}^{\infty} e^{-S(J)} \\
    &= \prod_{\alpha} \int_{-\infty}^{\infty} dx_\alpha \exp \Bigg [ \sum_{\alpha} \left ( -\frac{1}{2}m_\alpha x_\alpha^2 + J_\alpha x_\alpha \right ) \Bigg ] \\
    &= \prod_{\alpha} \left ( \int_{-\infty}^{\infty} dx_{\alpha} \exp \Bigg [ -\frac{1}{2}m_\alpha x_\alpha^2 + J_\alpha x_\alpha \Bigg ] \right ) \\
    &= \prod_{\alpha} \sqrt{\frac{2\pi}{m_\alpha}} e^{\frac{\frac{1}{2}J_\alpha\frac{1}{m_\alpha}J_\alpha}{2m_{\alpha}}} \\
    &= \frac{(2\pi)^{N/2}}{\sqrt{\det M}} e^{\frac{1}{2}\langle J | M^{-1}| J \rangle}
\end{align}


\section{Isserlis' Theorem}
Isserlis' Theorem is the statistical mechanical analogue of Wick's theorem and provides a way to compute expectations of products of multivariate normals. Let 
\begin{align}
(X_1,\dotsc, X_n) \in \mathcal{N}(\vec{0}, \Sigma)
\end{align}
be a zero mean multivariate normal. Then we can compute
\begin{align}
    E[X_1\dotsc X_n] &= \sum_{p \in P^2_n} \prod_{(i,j) \in p}  E[X_i X_j]
\end{align}
where $p$ is a pairing of $\{1,\dotsc, n\}$ into pairs and $P^2_n$ is the set of all such pairs.
\section{Feynman Diagrams}
Consider an expression of the form

\begin{align}
    \langle 0 | \mathcal{T}\{\phi_I(x_1)\dotsc \phi_I(x_n)\}|0 \rangle
\end{align}
We can apply Wick's theorem to the case where $n=4$ to represent it as the sum of three diagrams.
\begin{align}
\langle 0 | \mathcal{T} \{\phi_1\phi_2\phi_3\phi_4\}|0 \rangle &= 
\begin{tikzpicture}[baseline={([yshift=-.5ex]current bounding box.center)}]
\node[main] (1) {$1$}; 
\node[main] (2) [below of=1] {$3$};
\node[main] (3) [right of=1] {$2$};
\node[main] (4) [below of=3] {$4$};
\draw (1) -- (2);
\draw (3) -- (4);
\end{tikzpicture}
+
\begin{tikzpicture}[baseline={([yshift=-.5ex]current bounding box.center)}]
\node[main] (1) {$1$}; 
\node[main] (2) [right of=1] {$2$};
\node[main] (3) [below of=1] {$2$};
\node[main] (4) [below of=2] {$4$};
\draw (1) -- (2);
\draw (3) -- (4);
\end{tikzpicture}
+
\begin{tikzpicture}[baseline={([yshift=-.5ex]current bounding box.center)}]
\node[main] (1) {$1$}; 
\node[main] (2) [right of=1] {$2$};
\node[main] (3) [below of=1] {$2$};
\node[main] (4) [below of=2] {$4$};
\draw (1) -- (4);
\draw (3) -- (2);
\end{tikzpicture}
\end{align}
Physically, these expressions are usually interpreted as particles being created and annihilated at vertices. This intuition is often very useful, but we will also find it useful to view Feynman diagrams as purely combinatorial constructs as well.

\end{document}